\subsection*{一次完全准确的空查询}

重新取平台余额 10000、无收款记录，P1、C1 已在购买库提交。发送者发出 R1，R1 在网络中延迟。恢复者随后发查询 Q1，Q1 先到平台。平台此刻确实没有 C1 的效果，强查询返回空。Q1 的响应回到恢复者以后，R1 才到达并扣款，余额变为 7000。

\CSevenFigure{ch07-original05}{强查询返回空以后，原支付请求仍可到达}{fig:ch07-empty-inflight}

Q1 没有读错。将读取点放在 R1 到达以前，完全符合上一节的条件。空结果说明读取点尚无记录，R1 却还在途中。查询没有撤回它，也没有在平台写下禁止它执行的决定。

如果恢复者在 Q1 以后改用 \code{pay-C1-retry}，新请求 R2 可以先扣 3000，余 7000；R1 随后带 \code{pay-C1} 到达，再扣 3000，余 4000。平台对每个键只执行一次，仍符合自己的幂等规则。本地把一次 C1 分成两个下游身份，使两个请求分别取得执行资格。

若 R2 继续带 \code{pay-C1} 和原输入，平台让先到者完成，后到者读取同一收款。无论 R2 先扣、R1 晚到，还是 R1 先扣、R2 等待，都只产生一份效果。恢复者据此继续原 C1，无须猜测 R1 当前停在哪里。

沿新键分支核对，R2、R1 分别建立 X2、X1，两条流水各为 3000，最终余额为 4000。沿原键分支，两次到达找到同一个 X1，只有一条 3000 分流水，最终余额为 7000。两条分支中的空查询完全相同，效果数量取决于后来两个请求是否在平台争同一个执行位置。

\subsection*{只有查询时，恢复者少了什么}

假设平台接受商户关联 C1，也允许据此查询，却允许多个请求各带相同 C1，分别扣款。关联字段能帮助找记录，平台并未用它仲裁执行。查询不到后再发，即使仍带 C1，两个请求也可能各建一笔收款，把 10000 分扣到 4000 分。

将这个限制写成两条历史。\(H_{\text{丢失}}\) 中，R1 永久丢失；\(H_{\text{延迟}}\) 中，R1 留在途中，可以等待任意久以后再到达。在某个有限时刻以前，两条历史的每次强查询都可以返回空，本地也有相同的 P1、C1。为了让 \(H_{\text{丢失}}\) 自动推进，恢复者必须在有限步骤后决定再发一次。同样的观察会使它在 \(H_{\text{延迟}}\) 中作出同样决定，新请求和后来到达的 R1 就都有执行资格。

\CSevenFigure{ch07-original06}{仅查询时，永久丢失与任意延迟可以共享同一段观察}{fig:ch07-query-only}

这段推理依赖四项前提：请求可任意延迟，扣款可重复执行，查询只读当前记录，没有其他执行排他或关闭机制。查询可以完全准确。只要决定重发以前的观察相同，依赖这些观察的算法就会作出同样决定。随机选择重发，也会让某些允许历史具有重复的可能，达不到每条历史都满足效果上限的要求。

在这些前提下，可以只继续查询：R1 最后成功，就取得其证据；R1 永久丢失，购买保持未决。也可以选择再次发送来提高执行机会，同时明确承担重复扣款风险。要同时保证不重复扣款和自动进展，就要增加下游能力，或收紧请求延迟模型。

\subsection*{什么证据能够关闭旧执行机会}

一种补足方式是按固定身份去重，让 R1、R2 争同一个执行位置。另一种是提供原子关闭：平台关闭 C1 时，同时检查它是否已成功；已经成功就返回原收款，尚未成功则持久保存以后拒绝 C1 的决定。旧请求晚到，会读取关闭决定并被拒绝。查询到这样的关闭结果以后，服务才能根据业务规则创建下一动作。

关闭与扣款必须参与同一次仲裁。平台先查未付款，再用独立步骤写关闭，扣款可能从两步之间通过。只取消一条通信连接，也无法阻止另一条重复投递。所有执行路径都检查同一固定身份的关闭决定，才形成明确的旧执行边界。

另一类接口允许请求携带平台可核验的截止，由平台在真正取得执行资格时检查，拒绝过期旧请求。截止以前已受理的操作继续产生可查询的最终结局。只在客户端设超时，或只在购买库把 C1 标成取消，平台仍可能接受 R1。有限等待能够证明旧请求失效，依赖的是执行端截止、在途受理排空和查询完整性等具体条件。

取消支付对象还要解释：取消的是哪个对象，成功后是否仍允许捕获资金，以及它与并发付款怎样裁决。退款发生在已经收款之后，是带新身份的资金动作，保留原付款事实。先重复收款、发现后再退款，账本仍曾有两次扣款，用户也经历了资金占用和退款失败风险。

再沿一条错误历史走下去：Q1 强查空，R1 仍在途，服务把 P1 记为 \code{payment\_rejected}，告诉 Alice 付款失败，随后 R1 成功。平台余 7000，本地终态却声称付款拒绝。丢弃成功通知，会使两边长期不符；用成功覆盖已经发布的最终拒绝，又改变了拒绝结果的含义。准确的空查分支保持 \code{awaiting\_payment}，能够再次发送时继续原 C1，由后续证据决定终态。
